% Chapter V, Sections 51--52
% Chandrasekhar, "Hydrodynamic and Hydromagnetic Stability"
% FROM BELOW: 4. ROTATION AND MAGNETIC FIELD

\newcommand{\figurePlaceholder}[2]{%
  \begin{center}
    \fbox{\parbox{0.85\textwidth}{\centering\textbf{[Figure~#1]}\\[4pt]#2}}
  \end{center}
}

\section*{\S\,51. The perturbation equations}

% [Continues from previous page]
\noindent
\ldots\ last chapters, that the basic perturbation equations are (cf.\ equations
III~(78)--(80) and IV~(87)--(92)):
\begin{equation}
  \frac{\partial u_i}{\partial t}
  = -\frac{\partial}{\partial x_i}(\delta\varpi)
  + g\alpha\Theta\,\delta_{i3}
  + \nu\nabla^{2}u_i
  + 2\epsilon_{ijk}\,u_j\,\Omega_k
  + \frac{\mu H_j}{4\pi\rho}\,\frac{\partial h_i}{\partial x_j},
  \label{eq:5-21}
\end{equation}
\begin{equation}
  \frac{\partial h_i}{\partial t}
  = H_j\,\frac{\partial u_i}{\partial x_j}+\eta\nabla^{2}h_i,
  \label{eq:5-22}
\end{equation}
\begin{equation}
  \frac{\partial\Theta}{\partial t}
  = \beta w + \kappa\nabla^{2}\Theta,
  \label{eq:5-23}
\end{equation}
\begin{equation}
  \frac{\partial u_i}{\partial x_i} = 0, \quad\text{and}\quad
  \frac{\partial h_i}{\partial x_i} = 0.
  \label{eq:5-24}
\end{equation}

By taking the curl of equation~\eqref{eq:5-21}, we can eliminate the term
in~$\delta\varpi$; we thus obtain (cf.\ equations III~(82) and IV~(93))
\begin{equation}
  \frac{\partial\omega_i}{\partial t}
  = g\alpha\,\epsilon_{ijk}\,\frac{\partial\Theta}{\partial x_j}\,\lambda_k
  + \nu\nabla^{2}\omega_i
  + 2\Omega_j\,\frac{\partial u_i}{\partial x_j}
  + \frac{\mu H_j}{4\pi\rho}\,\frac{\partial v_i}{\partial x_j},
  \label{eq:5-25}
\end{equation}
where
\begin{equation}
  \boldsymbol{\omega} = \operatorname{curl}\mathbf{u}
  \quad\text{and}\quad
  \mathbf{v} = \operatorname{curl}\mathbf{h}.
  \label{eq:5-26}
\end{equation}

Taking the curl of equation~\eqref{eq:5-25} once again, we obtain
(cf.\ equations III~(83) and IV~(95))
\begin{equation}
  \frac{\partial}{\partial t}\nabla^{2}u_i
  = g\alpha\!\left(
      \lambda_i\,\nabla^{2}\Theta
      - \lambda_j\,\frac{\partial^{2}\Theta}{\partial x_i\,\partial x_j}
    \right)
  + \nu\nabla^{4}u_i
  - 2\Omega_j\,\frac{\partial\omega_i}{\partial x_j}
  + \frac{\mu H_j}{4\pi\rho}\,\frac{\partial}{\partial x_j}\,\nabla^{2}h_i.
  \label{eq:5-27}
\end{equation}

Now multiplying equations~\eqref{eq:5-25} and~\eqref{eq:5-27} by~$\delta_{i3}$
we get
\begin{equation}
  \frac{\partial\zeta}{\partial t}
  = \nu\nabla^{2}\zeta
  + 2\Omega\,\frac{\partial w}{\partial z}
  + \frac{\mu H}{4\pi\rho}\,\frac{\partial\xi}{\partial z},
  \label{eq:5-28}
\end{equation}
and
\begin{equation}
  \frac{\partial}{\partial t}\nabla^{2}w
  = g\alpha\!\left(
      \frac{\partial^{2}\Theta}{\partial x^{2}}
      + \frac{\partial^{2}\Theta}{\partial y^{2}}
    \right)
  + \nu\nabla^{4}w
  - 2\Omega\,\frac{\partial\zeta}{\partial z}
  + \frac{\mu H}{4\pi\rho}\,\frac{\partial}{\partial z}\,\nabla^{2}h_3,
  \label{eq:5-29}
\end{equation}
where $\zeta$ and~$\xi/4\pi$ are the $z$-components of the vorticity and the
current density, respectively.  Equations~\eqref{eq:5-28} and~\eqref{eq:5-29}
replace equations~IV~(99) and~(100); the remaining equations IV~(96)--(98) are
unaffected.

We shall restrict our discussion of this problem to the case when
$\mathbf{g}$, $\mathbf{H}$, and~$\boldsymbol{\Omega}$ act in the same
direction.  In this case the relevant equations are:
\begin{equation}
  \frac{\partial\Theta}{\partial t}
  = \kappa\nabla^{2}\Theta + \beta w,
  \label{eq:5-30}
\end{equation}
\begin{equation}
  \frac{\partial h_3}{\partial t}
  = \eta\nabla^{2}h_3 + H\,\frac{\partial w}{\partial z},
  \label{eq:5-31}
\end{equation}
\begin{equation}
  \frac{\partial\xi}{\partial t}
  = \eta\nabla^{2}\xi + H\,\frac{\partial\zeta}{\partial z},
  \label{eq:5-32}
\end{equation}
\begin{equation}
  \frac{\partial\zeta}{\partial t}
  = \nu\nabla^{2}\zeta
  + 2\Omega\,\frac{\partial w}{\partial z}
  + \frac{\mu H}{4\pi\rho}\,\frac{\partial\xi}{\partial z},
  \label{eq:5-33}
\end{equation}
and
\begin{equation}
  \frac{\partial}{\partial t}(\nabla^{2}w)
  = g\alpha\!\left(
      \frac{\partial^{2}\Theta}{\partial x^{2}}
      + \frac{\partial^{2}\Theta}{\partial y^{2}}
    \right)
  + \nu\nabla^{4}w
  - 2\Omega\,\frac{\partial\zeta}{\partial z}
  + \frac{\mu H}{4\pi\rho}\,\frac{\partial}{\partial z}\,\nabla^{2}h_3.
  \label{eq:5-34}
\end{equation}

Analysing the disturbances into two-dimensional waves, and considering
disturbances characterized by a particular wave number, we find that
equations~\eqref{eq:5-30}--\eqref{eq:5-34} give
\begin{equation}
  (D^{2}-a^{2}-p_1\,\sigma)\,\Theta
  = -\left(\frac{\beta d^{2}}{\kappa}\right)W,
  \label{eq:5-35}
\end{equation}
\begin{equation}
  (D^{2}-a^{2}-p_2\,\sigma)\,K
  = -\left(\frac{Hd}{\eta}\right)DW,
  \label{eq:5-36}
\end{equation}
\begin{equation}
  (D^{2}-a^{2}-p_2\,\sigma)\,X
  = -\left(\frac{Hd}{\eta}\right)DZ,
  \label{eq:5-37}
\end{equation}
\begin{equation}
  (D^{2}-a^{2}-\sigma)\,Z
  = -\left(\frac{2\Omega d^{2}}{\nu}\right)DW
  - \left(\frac{\mu Hd}{4\pi\rho\nu}\right)DX,
  \label{eq:5-38}
\end{equation}
and
\begin{multline}
  (D^{2}-a^{2})
  (D^{2}-a^{2}-\sigma)\,W
  + \left(\frac{\mu Hd}{4\pi\rho\nu}\right)
  D(D^{2}-a^{2})\,K \\
  - \left(\frac{2\Omega d^{2}}{\nu}\right)DZ
  = \left(\frac{g\alpha d^{2}}{\nu}\right)a^{2}\Theta,
  \label{eq:5-39}
\end{multline}
where the notation is the same as in equations~IV~(118)--(123).

Eliminating $X$ between equations~\eqref{eq:5-37} and~\eqref{eq:5-38},
and similarly~$K$ between equations~\eqref{eq:5-36}
and~\eqref{eq:5-39}, we obtain
\begin{multline}
  \bigl[(D^{2}-a^{2}-\sigma)(D^{2}-a^{2}-p_2\,\sigma)
    -QD^{2}\bigr]Z \\
  = -\left(\frac{2\Omega d^{2}}{\nu}\right)
  D(D^{2}-a^{2}-p_2\,\sigma)\,W
  \label{eq:5-40}
\end{multline}
and
\begin{multline}
  \bigl[D^{2}-a^{2}\bigr]
  \bigl[(D^{2}-a^{2}-\sigma)(D^{2}-a^{2}-p_2\,\sigma)
    -QD^{2}\bigr]W \\
  -\left(\frac{2\Omega d^{2}}{\nu}\right)
  D(D^{2}-a^{2}-p_2\,\sigma)\,Z
  = \left(\frac{g\alpha d^{2}}{\nu}\right)
  a^{2}(D^{2}-a^{2}-p_2\,\sigma)\,\Theta.
  \label{eq:5-41}
\end{multline}

Eliminating $Z$ between these last two equations, we have
\begin{multline}
  \bigl\{(D^{2}\!-\!a^{2})
    \bigl[(D^{2}\!-\!a^{2}\!-\!\sigma)(D^{2}\!-\!a^{2}\!-\!p_2\,\sigma)
      -QD^{2}\bigr]^{2} \\
  +\,T D^{2}(D^{2}\!-\!a^{2}\!-\!p_2\,\sigma)^{2}\bigr\}W \\
  = \left(\frac{g\alpha d^{2}}{\nu}\right)
  a^{2}\bigl[(D^{2}\!-\!a^{2}\!-\!\sigma)(D^{2}\!-\!a^{2}\!-\!p_2\,\sigma)
    -QD^{2}\bigr]
  (D^{2}\!-\!a^{2}\!-\!p_2\,\sigma)\,\Theta.
  \label{eq:5-42}
\end{multline}

Finally, eliminating~$\Theta$ between equations~\eqref{eq:5-35}
and~\eqref{eq:5-42} we obtain
\begin{multline}
  (D^{2}\!-\!a^{2}\!-\!p_1\,\sigma)
  \bigl\{(D^{2}\!-\!a^{2})
    \bigl[(D^{2}\!-\!a^{2}\!-\!\sigma)(D^{2}\!-\!a^{2}\!-\!p_2\,\sigma)
      -QD^{2}\bigr]^{2} \\
  +\,T D^{2}(D^{2}\!-\!a^{2}\!-\!p_2\,\sigma)^{2}\bigr\}W \\
  = -\mathit{Ra}^{2}
    \bigl[(D^{2}\!-\!a^{2}\!-\!\sigma)(D^{2}\!-\!a^{2}\!-\!p_2\,\sigma)
    -QD^{2}\bigr]
    (D^{2}\!-\!a^{2}\!-\!p_2\,\sigma)\,W.
  \label{eq:5-43}
\end{multline}

Solutions of the foregoing equations must be sought which satisfy the
boundary conditions given in equations~II~(95) and~(96) and IV~(124)
and~(125).

It may be recalled here that all equations following~\eqref{eq:5-21} apply only if
$\mathbf{g}$, $\mathbf{H}$, and~$\boldsymbol{\Omega}$ are parallel.  The
generalization to the case when $\mathbf{g}$, $\mathbf{H}$,
and~$\boldsymbol{\Omega}$ are coplanar is immediate.  By restricting oneself
to the onset of convection in the form of infinitely extended rolls (along
directions parallel to the plane containing $\mathbf{g}$, $\mathbf{H}$,
and~$\boldsymbol{\Omega}$), instead of cells, one can readily verify that the
same set of equations applies to this more general case if we interpret $H$
and~$\Omega$ in all the equations to mean the components of~$\mathbf{H}$
and~$\boldsymbol{\Omega}$ in the direction of the vertical.  The restriction
to the onset of convection in the form of rolls when $\mathbf{g}$,
$\mathbf{H}$, and~$\boldsymbol{\Omega}$ are coplanar can very likely be
justified in the way in which the same restriction was justified in
Chapter~IV, \S\,47 for the case when $\mathbf{H}$ alone was acting.
However, when $\mathbf{g}$, $\mathbf{H}$, and~$\boldsymbol{\Omega}$ are not
coplanar, the solution of the problem will require an explicit minimizing of
the Rayleigh number as a function of the two wave numbers characterizing the
disturbance in two directions at right angles to one another and in the
horizontal plane.  But it is unlikely that the solution will exhibit any
essentially new feature not already disclosed by the solution for the case
when $\mathbf{g}$, $\mathbf{H}$, and~$\boldsymbol{\Omega}$ are parallel.

\section*{\S\,52. The case when instability sets in as stationary convection}

We shall first consider the case when instability sets in as ordinary
convection.  The case when it sets in as overstability is considered in the
following section.

When instability sets in as ordinary convection, the marginal state will be
characterized by $\sigma = 0$; and the basic equations are:
\begin{equation}
  (D^{2}-a^{2})\,\Theta
  = -\left(\frac{\beta d^{2}}{\kappa}\right)W,
  \label{eq:5-44}
\end{equation}
\begin{equation}
  \bigl[(D^{2}-a^{2})^{2}-QD^{2}\bigr]Z
  = -\left(\frac{2\Omega d}{\nu}\right)D(D^{2}-a^{2})\,W,
  \label{eq:5-45}
\end{equation}
\begin{equation}
  (D^{2}-a^{2})\,K
  = -\left(\frac{Hd}{\eta}\right)DW,
  \label{eq:5-46}
\end{equation}
\begin{equation}
  (D^{2}-a^{2})\,X
  = -\left(\frac{Hd}{\eta}\right)DZ,
  \label{eq:5-47}
\end{equation}
and
\begin{multline}
  (D^{2}-a^{2})
  \bigl\{\bigl[(D^{2}-a^{2})^{2}-QD^{2}\bigr]^{2}
  +T D^{2}(D^{2}-a^{2})^{2}\bigr\}W \\
  = -\mathit{Ra}^{2}\bigl\{(D^{2}-a^{2})^{2}-QD^{2}\bigr\}W.
  \label{eq:5-48}
\end{multline}

According to equations~II~(95) and~(96) and IV~(124) and~(125), the
boundary conditions with respect to which equations~\eqref{eq:5-44}--\eqref{eq:5-48}
must be solved are:
\begin{equation}
  W = 0 \quad\text{and}\quad \Theta = 0 \quad\text{for } z = 0 \text{ and } 1,
  \label{eq:5-49}
\end{equation}
\begin{equation}
  \left.
  \begin{aligned}
    &\textit{either}\quad DW = 0 \quad\text{and}\quad Z = 0
      \quad\text{(on a rigid boundary)} \\
    &\textit{or}\quad D^{2}W = 0 \quad\text{and}\quad DZ = 0
      \quad\text{(on a free boundary)}
  \end{aligned}
  \right\}
  \label{eq:5-50}
\end{equation}
and
\begin{equation}
  \left.
  \begin{aligned}
    &\textit{either}\quad DX = 0 \;\text{and}\; K = 0
      \;\text{(on a perfectly conducting boundary)}, \\
    &\textit{or}\quad X = 0
      \;\text{(on a boundary adjoining a non-conducting medium)}.
  \end{aligned}
  \right\}
  \label{eq:5-51}
\end{equation}

The foregoing equations and boundary conditions constitute a characteristic
value problem in a differential equation of order ten.  In spite of a
variational formulation which is possible (see Appendix~II), the problem is
clearly of considerable complexity.  For this reason, we shall restrict
ourselves to the case when both the bounding surfaces are free and the medium
adjoining the fluid is non-conducting.  While this is admittedly an artificial
case to consider, we do not expect to lose any of the essential features of the
problem on this account; for, we have seen in Chapters~III and~IV that the
solutions obtained for different sets of boundary conditions all exhibit the
same dependence on the parameters of the problem such as $\mathrm{p}$, $T$,
and~$Q$.

\subsection*{(a) \textit{The solution for the case of two free boundaries}}

As in Chapter~III, \S\,27\,(a) and Chapter~IV, \S\,44\,(a), it can be shown
that for this case, the proper solutions of
equations~\eqref{eq:5-44}--\eqref{eq:5-48} are given by
\begin{equation}
  W = W_0\sin n\pi z,
  \label{eq:5-52}
\end{equation}
\begin{equation}
  \Theta = \frac{\beta}{\kappa}\,
  \frac{d^{2}}{n^{2}\pi^{2}+a^{2}}\,W_0\sin n\pi z,
  \label{eq:5-53}
\end{equation}
\begin{equation}
  Z = \frac{2\Omega d}{\nu}\,
  \frac{n\pi(n^{2}\pi^{2}+a^{2})}{(n^{2}\pi^{2}+a^{2})^{2}+Qn^{2}\pi^{2}}\,
  W_0\cos n\pi z,
  \label{eq:5-54}
\end{equation}
\begin{equation}
  K = \frac{Hd}{\eta}\,
  \frac{n\pi}{n^{2}\pi^{2}+a^{2}}\,
  W_0\cos n\pi z + \text{constant} \times \cosh az,
  \label{eq:5-55}
\end{equation}
and
\begin{equation}
  X = \frac{2\Omega d}{\nu}\,\frac{Hd}{\eta}\,
  \frac{n^{2}\pi^{2}}{(n^{2}\pi^{2}+a^{2})^{2}+Q n^{2}\pi^{2}}\,
  W_0\sin n\pi z,
  \label{eq:5-56}
\end{equation}
where $n$ is an integer.  The corresponding characteristic equation is
obtained by substituting the solution~\eqref{eq:5-52} for~$W$ in
equation~\eqref{eq:5-48}.  We find
\begin{equation}
  R = \frac{%
    (n^{2}\pi^{2}+a^{2})
    \bigl\{\bigl[(n^{2}\pi^{2}+a^{2})^{2}+Qn^{2}\pi^{2}\bigr]^{2}
    +Tn^{2}\pi^{2}(n^{2}\pi^{2}+a^{2})\bigr\}}{%
    a^{2}\bigl[(n^{2}\pi^{2}+a^{2})^{2}+Qn^{2}\pi^{2}\bigr]}.
  \label{eq:5-57}
\end{equation}

Letting $a^{2} = \pi^{2}x$, we can rewrite equation~\eqref{eq:5-57} in
the form
\begin{equation}
  R = \pi^{4}\,\frac{n^{2}+x}{x}
  \left[
    (n^{2}+x)^{2}
    + \frac{Qn^{2}}{\pi^{2}}
  \right]
  + \frac{T}{x}\,
  \frac{1}{1/n^{2}+Q/\pi^{2}(n^{2}+x)^{2}}.
  \label{eq:5-58}
\end{equation}

From this equation it follows that instability first sets in for the lowest
mode $n = 1$.  The corresponding expression for~$R$ is
\begin{equation}
  R = \pi^{4}\,\frac{%
    (1+x)
    \bigl\{\bigl[(1+x)^{2}+Q_1\bigr]^{2}
    +T_1(1+x)\bigr\}}{%
    x\bigl[(1+x)^{2}+Q_1\bigr]},
  \label{eq:5-59}
\end{equation}
where
\begin{equation}
  x = \frac{a^{2}}{\pi^{2}}, \quad
  Q_1 = \frac{Q}{\pi^{2}}, \quad\text{and}\quad
  T_1 = \frac{T}{\pi^{4}}.
  \label{eq:5-60}
\end{equation}

As a function of~$x$, $R$ given by equation~\eqref{eq:5-59} attains its
extremal value when
\begin{equation}
  2x^{3}+3x^{2}-1
  = Q_1 + T_1\,
  \frac{(1+x)^{4}-Q_1(x^{2}-1)}{%
    \bigl[(1+x)^{2}+Q_1\bigr]^{2}}.
  \label{eq:5-61}
\end{equation}

\figurePlaceholder{47}{The critical Rayleigh number $R_c$ for the cases of
ordinary cellular convection (solid line) and overstability (for
$\mathrm{p} = 0.025$) (broken line) as a function of
$Q_1\,(= Q/\pi^{2})$ for various assigned values of
$T_1\,(= T/\pi^{4})$.  The curves are labeled by the values of~$T_1$ to which
they refer.  For a given value of~$T_1$, instability will set in as
overstability for all values of~$Q_1$ less than that at the point of
intersection of the corresponding full-line and dashed curves; for all larger
values of~$Q_1$, it will set in as ordinary convection.}

However, this equation is not very useful for determining the critical Rayleigh
numbers for assigned values of~$Q_1$ and~$T_1$.  It is more convenient to
evaluate~$R$ directly as a function of~$x$ (in accordance with
equation~\eqref{eq:5-59}) and locate the minimum numerically.  The critical
numbers listed in Table~XVIII and illustrated in Figs.~47 and~48 were
determined in this fashion.

It is apparent from Table~XVIII and Figs.~47 and~48 that the solution

% ----------------------------------------------------------------
% TABLE XVIII
% ----------------------------------------------------------------
\begin{table}[p]
\centering
\caption{The critical Rayleigh numbers and the wave numbers of the associated
disturbance for the onset of instability as stationary convection for various
values of $Q_1\,(= Q/\pi^{2})$ and $T_1\,(= T/\pi^{4})$.}
\label{tab:XVIII}

\medskip
\small
\begin{tabular}{r@{\quad}r@{\quad}r|r@{\quad}r@{\quad}r}
\hline
\multicolumn{3}{c|}{$T_1 = 0$} & \multicolumn{3}{c}{$T_1 = 10$} \\
$Q_1$ & $a$ & $R_c$ & $Q_1$ & $a$ & $R_c$ \\
\hline
 0     & $3{\cdot}70$  &           & 0     & $4{\cdot}31$  &           \\
 10    & $3{\cdot}97$  & $6{\cdot}95\times10^{3}$ & 10    & $4{\cdot}56$  &           \\
 40    &               &           & 40    &               &           \\
 100   & $5{\cdot}67$  &           & 100   & $5{\cdot}96$  &           \\
 1{,}000 & $8{\cdot}59$ &          & 1{,}000 & $8{\cdot}65$ &          \\
 10{,}000 & $12{\cdot}6$ &$1{\cdot}06\times10^{6}$& 10{,}000 & $12{\cdot}6$  & \\
 50{,}000 & $16{\cdot}8$ &$5{\cdot}09\times10^{6}$& 50{,}000 & $16{\cdot}8$  & \\
 100{,}000 & $18{\cdot}94$&$1{\cdot}01\times10^{7}$& 100{,}000 & $18{\cdot}94$ & \\
\hline
\end{tabular}

\bigskip
\begin{tabular}{r@{\quad}r@{\quad}r|r@{\quad}r@{\quad}r}
\hline
\multicolumn{3}{c|}{$T_1 = 50$} & \multicolumn{3}{c}{$T_1 = 100$} \\
$Q_1$ & $a$ & $R_c$ & $Q_1$ & $a$ & $R_c$ \\
\hline
 0    & $5{\cdot}04$  & $3{\cdot}88\times10^{4}$ & 0    & $5{\cdot}73$ &$1{\cdot}03\times10^{5}$\\
 10   & $5{\cdot}30$  &                          & 10   & $5{\cdot}92$ &           \\
 40   & $5{\cdot}96$  & $6{\cdot}43\times10^{4}$ & 40   & $6{\cdot}45$ &           \\
 100  & $6{\cdot}95$  &                          & 100  & $7{\cdot}20$ &           \\
 1{,}000 & $10{\cdot}13$ & $1{\cdot}14\times10^{5}$& 1{,}000 &             &           \\
 10{,}000 & $12{\cdot}6$ &                        & 10{,}000 & $12{\cdot}6$ &          \\
 50{,}000 & $16{\cdot}8$ &                        & 50{,}000 & $16{\cdot}8$ &          \\
 100{,}000 & $18{\cdot}94$ &                      & 100{,}000 & $18{\cdot}94$ &        \\
\hline
\end{tabular}

\bigskip
\begin{tabular}{r@{\quad}r@{\quad}r|r@{\quad}r@{\quad}r}
\hline
\multicolumn{3}{c|}{$T_1 = 500$} & \multicolumn{3}{c}{$T_1 = 1{,}000$} \\
$Q_1$ & $a$ & $R_c$ & $Q_1$ & $a$ & $R_c$ \\
\hline
 0    & $7{\cdot}90$  & $5{\cdot}03\times10^{5}$ & 0    & $9{\cdot}05$  & $2{\cdot}05\times10^{6}$\\
 10   & $7{\cdot}98$  &                          & 20   & $9{\cdot}14$  &           \\
 20   & $8{\cdot}08$  & $1{\cdot}73\times10^{5}$ & 40   & $9{\cdot}28$  &           \\
 40   & $8{\cdot}39$  &                          & 50   & $9{\cdot}36$  &           \\
 50   & $8{\cdot}60$  & $2{\cdot}05\times10^{5}$ & 80   & $9{\cdot}60$  &           \\
 80   & $9{\cdot}19$  &                          & 100  & $9{\cdot}80$  &           \\
 100  & $9{\cdot}52$  & $3{\cdot}45\times10^{5}$ & 500  & $11{\cdot}78$ &           \\
 500  & $11{\cdot}46$ & $6{\cdot}93\times10^{5}$ & 1{,}000 & $12{\cdot}94$ &        \\
 1{,}000 & $12{\cdot}77$ & $1{\cdot}08\times10^{6}$& 10{,}000 & $18{\cdot}0$ &      \\
 10{,}000 & $17{\cdot}8$ & $8{\cdot}65\times10^{6}$& 50{,}000 &             &       \\
 50{,}000 & $16{\cdot}8$ &                        & 100{,}000 & $18{\cdot}94$ &      \\
 100{,}000 & $18{\cdot}94$ & $1{\cdot}01\times10^{7}$ & & & \\
\hline
\end{tabular}

\bigskip
\begin{tabular}{r@{\quad}r@{\quad}r|r@{\quad}r@{\quad}r}
\hline
\multicolumn{3}{c|}{$T_1 = 1{,}500$} & \multicolumn{3}{c}{$T_1 = 2{,}000$} \\
$Q_1$ & $a$ & $R_c$ & $Q_1$ & $a$ & $R_c$ \\
\hline
 0    & $9{\cdot}87$  & $4{\cdot}05\times10^{6}$ & 0    & $10{\cdot}5$ & $6{\cdot}78\times10^{6}$\\
 10   & $9{\cdot}90$  &                          & 10   & $10{\cdot}5$ &           \\
 20   & $9{\cdot}95$  &                          & 20   & $10{\cdot}5$ &           \\
 30   & $10{\cdot}0$  & $4{\cdot}53\times10^{6}$ & 30   & $10{\cdot}5$ &           \\
 40   & $10{\cdot}1$  &                          & 40   & $10{\cdot}6$ &           \\
 50   & $10{\cdot}1$  & $5{\cdot}14\times10^{6}$ & 50   & $10{\cdot}6$ & $8{\cdot}17\times10^{6}$\\
 80   & $10{\cdot}4$  &                          & 80   & $10{\cdot}8$ &           \\
 100  & $10{\cdot}6$  & $6{\cdot}55\times10^{6}$ & 100  & $10{\cdot}9$ &           \\
 500  & $12{\cdot}20$ &                          & 500  & $12{\cdot}57$ &          \\
 1{,}000 & $13{\cdot}25$ &                       & 1{,}000 & $13{\cdot}60$ &       \\
 10{,}000 & $18{\cdot}2$ &$8{\cdot}65\times10^{6}$& 10{,}000 & $18{\cdot}3$ &     \\
 50{,}000 & $16{\cdot}94$ &                      & 50{,}000 &              &       \\
 100{,}000 & $18{\cdot}94$ & $1{\cdot}02\times10^{7}$ & 100{,}000 & $18{\cdot}94$ & \\
\hline
\end{tabular}

\bigskip
\begin{tabular}{r@{\quad}r@{\quad}r|r@{\quad}r@{\quad}r}
\hline
\multicolumn{3}{c|}{$T_1 = 5{,}000$} & \multicolumn{3}{c}{$T_1 = 10{,}000$} \\
$Q_1$ & $a$ & $R_c$ & $Q_1$ & $a$ & $R_c$ \\
\hline
 10   & $12{\cdot}02$ & $4{\cdot}00\times10^{7}$ & 10   & $12{\cdot}98$ & $6{\cdot}09\times10^{7}$\\
 20   & $9{\cdot}95$  &                          & 20   &              & $8{\cdot}88\times10^{7}$\\
 40   & $9{\cdot}82$  & $4{\cdot}00\times10^{7}$ & 40   & $11{\cdot}09$ &           \\
 50   & $7{\cdot}00$  & $3{\cdot}93\times10^{7}$ & 50   & $8{\cdot}69$  &           \\
 80   & $4{\cdot}94$  &                          & 80   & $8{\cdot}68$  &           \\
 100  & $4{\cdot}55$  & $3{\cdot}77\times10^{7}$ & 100  & $3{\cdot}77$  &           \\
 200  & $4{\cdot}13$  &                          & 200  & $4{\cdot}73$  & $7{\cdot}21\times10^{7}$\\
 500  & $5{\cdot}66$  &                          & 500  & $5{\cdot}31$  &           \\
 1{,}000 & $7{\cdot}22$ & $6{\cdot}89\times10^{7}$& 1{,}000 & $6{\cdot}51$ &       \\
 10{,}000 & $12{\cdot}6$ & $1{\cdot}21\times10^{8}$& 10{,}000 &            &       \\
 50{,}000 &             &                        & 50{,}000 &              &       \\
 100{,}000 & $18{\cdot}94$ &$1{\cdot}03\times10^{7}$& 100{,}000 & $18{\cdot}94$ & \\
\hline
\end{tabular}

\bigskip
\begin{tabular}{r@{\quad}r@{\quad}r|r@{\quad}r@{\quad}r}
\hline
\multicolumn{3}{c|}{$T_1 = 50{,}000$} & \multicolumn{3}{c}{$T_1 = 100{,}000$} \\
$Q_1$ & $a$ & $R_c$ & $Q_1$ & $a$ & $R_c$ \\
\hline
 10   & $19{\cdot}3$ & $1{\cdot}86\times10^{8}$ & 10   & $18{\cdot}9$ & $4{\cdot}15\times10^{8}$\\
 20   & $19{\cdot}1$ &                          & 20   & $18{\cdot}8$ &           \\
 40   & $19{\cdot}1$ & $1{\cdot}81\times10^{8}$ & 40   & $18{\cdot}4$ &           \\
 50   & $14{\cdot}9$  & $1{\cdot}86\times10^{8}$& 50   & $18{\cdot}2$ & $3{\cdot}96\times10^{8}$\\
 80   & $14{\cdot}8$  &                         & 80   &              &           \\
 100  & $14{\cdot}54$ & $3{\cdot}86\times10^{8}$& 100  & $3{\cdot}54$ & $1{\cdot}37\times10^{8}$\\
 200  & $3{\cdot}97$ & $5{\cdot}94\times10^{7}$ & 200  & $3{\cdot}48$ &           \\
 500  & $3{\cdot}71$  &                         & 500  & $4{\cdot}33$ &           \\
 1{,}000 & $3{\cdot}09$ &                       & 1{,}000 & $3{\cdot}97$ & $1{\cdot}37\times10^{8}$\\
 4{,}000 &             &                        & 4{,}000 &              &           \\
 10{,}000 & $10{\cdot}5$ &$1{\cdot}08\times10^{8}$& 10{,}000 &            &          \\
 30{,}000 &             &                        & 30{,}000 &              &         \\
 50{,}000 & $16{\cdot}6$ &                       & 50{,}000 &              &         \\
 100{,}000 & $18{\cdot}94$ & $1{\cdot}03\times10^{7}$ & 100{,}000 & $18{\cdot}94$ & \\
\hline
\end{tabular}

\bigskip
\begin{tabular}{r@{\quad}r@{\quad}r|r@{\quad}r@{\quad}r}
\hline
\multicolumn{3}{c|}{$T_1 = 10^{5}$} & \multicolumn{3}{c}{$T_1 = 10^{6}$} \\
$Q_1$ & $a$ & $R_c$ & $Q_1$ & $a$ & $R_c$ \\
\hline
 10   & $47{\cdot}8$ & $1{\cdot}85\times10^{8}$ & 10   & $41{\cdot}01$ &           \\
 50   & $47{\cdot}2$ &                          & 50   & $41{\cdot}01$ &           \\
 80   & $47{\cdot}2$ &                          & 80   & $35{\cdot}43$ &           \\
 100  & $47{\cdot}3$ & $1{\cdot}44\times10^{8}$ & 100  &              &           \\
 200  & $37{\cdot}1$ &                          & 200  &              &           \\
 300  & $3{\cdot}35$ &                          & 300  & $3{\cdot}19$ &           \\
 500  & $3{\cdot}12$ &                          & 500  & $3{\cdot}10$ &           \\
 1{,}000 & $3{\cdot}35$ &$9{\cdot}03\times10^{7}$& 1{,}000 &            &          \\
 5{,}000 & $5{\cdot}13$ &                       & 5{,}000 &              &         \\
 10{,}000 & $9{\cdot}33$ &$1{\cdot}06\times10^{8}$& 10{,}000 &           &          \\
 100{,}000 & $17{\cdot}9$ &                     & 100{,}000 &             &          \\
 1{,}000{,}000 & $18{\cdot}94$ &                & 1{,}000{,}000 & $18{\cdot}94$ &   \\
\hline
\end{tabular}

\bigskip
\begin{tabular}{r@{\quad}r@{\quad}r|r@{\quad}r@{\quad}r}
\hline
\multicolumn{3}{c|}{$T_1 = 10^{7}$} & \multicolumn{3}{c}{$T_1 = 10^{8}$} \\
$Q_1$ & $a$ & $R_c$ & $Q_1$ & $a$ & $R_c$ \\
\hline
 10   & $60{\cdot}96$ & $3{\cdot}97\times10^{9}$ & 50   & $88{\cdot}47$ &           \\
 50   & $60{\cdot}23$ &                          & 100  & $88{\cdot}45$ &           \\
 100  & $60{\cdot}14$ & $3{\cdot}97\times10^{9}$ & 200  & $88{\cdot}43$ &           \\
 500  & $60{\cdot}86$ & $1{\cdot}38\times10^{9}$ & 500  & $88{\cdot}34$ &$1{\cdot}93\times10^{9}$\\
 800  & $60{\cdot}86$ & $4{\cdot}86\times10^{8}$ & 1{,}000 & $88{\cdot}24$ &       \\
 1{,}000 & $3{\cdot}47$ &                       & 3{,}000 & $3{\cdot}96$ &          \\
 5{,}000 & $3{\cdot}14$ &$3{\cdot}96\times10^{8}$& 5{,}000 & $3{\cdot}15$ &        \\
 10{,}000 & $3{\cdot}74$ &                      & 10{,}000 & $3{\cdot}36$ &$1{\cdot}03\times10^{8}$\\
 100{,}000 & $17{\cdot}4$ & $1{\cdot}03\times10^{8}$ & 100{,}000 &         &        \\
 1{,}000{,}000 & $18{\cdot}94$ &                & 1{,}000{,}000 &$18{\cdot}94$&     \\
\hline
\end{tabular}

\bigskip
\begin{tabular}{r@{\quad}r@{\quad}r}
\hline
\multicolumn{3}{c}{$T_1 = 10^{9}$} \\
$Q_1$ & $a$ & $R_c$ \\
\hline
 50   & $130{\cdot}2$  & $8{\cdot}35\times10^{9}$ \\
 100  & $130{\cdot}2$  &           \\
 500  & $130{\cdot}9$  & $8{\cdot}15\times10^{9}$ \\
 1{,}000 & $130{\cdot}2$  &       \\
 2{,}000 & $130{\cdot}9$  & $8{\cdot}54\times10^{9}$ \\
 5{,}000 & $130{\cdot}9$ & $7{\cdot}79\times10^{9}$ \\
 10{,}000 & $3{\cdot}14$  & $3{\cdot}56\times10^{8}$\\
 50{,}000 & $3{\cdot}14$  &       \\
 100{,}000 & $3{\cdot}71$ & $5{\cdot}07\times10^{8}$ \\
 300{,}000 & $7{\cdot}47$ & $3{\cdot}43\times10^{8}$ \\
 1{,}000{,}000 & $17{\cdot}2$  &  \\
 3{,}000{,}000 & $18{\cdot}94$ & $3{\cdot}43\times10^{8}$ \\
 10{,}000{,}000 &             &   \\
 30{,}000{,}000 & $44{\cdot}6$ & $8{\cdot}94\times10^{8}$ \\
 300{,}000{,}000 &              & $8{\cdot}90$  \\
\hline
\end{tabular}
\end{table}

% ----------------------------------------------------------------
% Figures 47 and 48
% ----------------------------------------------------------------

\figurePlaceholder{48}{The dependence on $Q_1$ (for various assigned values
of~$T_1$) of the wave number~$a$ (in the unit $1/d$) of the disturbance at
which instability first sets in as convection (solid line) and as
overstability (for $\mathrm{p} = 0{\cdot}025$) (broken line).  It will be
observed that a discontinuous change in~$a$ occurs when (for increasing~$Q_1$)
the manner of instability changes from overstability to cellular convection.}

of the problem presents some very unexpected features.  These result from
the fact that the curve $R(a)$ defined by equations~\eqref{eq:5-59}
and~\eqref{eq:5-60} has, for certain ranges of the parameters~$Q_1$
and~$T_1$, \emph{two minima} (see Fig.~49); and the minimum giving the lower
Rayleigh number, when $Q_1$ is less than a certain value, gives the higher
Rayleigh number when $Q_1$ is larger than this value.  Thus, for
$T_1 = 10^{4}$ and $Q_1 = 80$, the two minima occur

\figurePlaceholder{49}{The dependence of the Rayleigh number~$R_c$ for the
onset of instability on the wave number when $T_1 = 10^{4}$ and $Q_1$ has the
values 40 (curve labelled~a), 100 (curve labelled~b), and 200 (curve
labelled~c).  The abscissa is related to the wave number~$a$ by
$x = a^{2}/\pi^{2}$.  Notice the occurrence of two minima and their relative
disposition.}

for $a = 18{\cdot}3$ and $3{\cdot}38$ where
$R = 4{\cdot}00\times10^{4}$ and $4{\cdot}78\times10^{4}$, respectively;
while, for $Q_1 = 100$ the two minima occur for $a = 18{\cdot}2$ and
$3{\cdot}37$, where $R = 3{\cdot}98\times10^{4}$ and
$3{\cdot}93\times10^{4}$, respectively.  Consequently, for $Q_1$ slightly
less than~100, the wave number of the cells which appear at marginal
stability will suddenly decrease from $a = 18{\cdot}2$ to $a = 3{\cdot}4$.
In other words, if we start with an initial situation in which
$T_1 = 10^{4}$ and no magnetic field is present, and gradually increase the
strength of the magnetic field, then at first the cells which appear at
marginal stability will be elongated; and when the magnetic field has
increased to a value corresponding to $Q_1 = 100$, cells of two very
different sizes will appear simultaneously: one set which will be highly
elongated and another set which will, relatively, be highly flattened.  As the
magnetic field increases
